% Source: 60_Chapter_46_Christianity_In_the_West_Decline_and_Re.txt
\bookchapter[Christianity in the West]{CHAPTER 46}{CHRISTIANITY IN THE WEST}{Decline and Reconstruction}

\leadin{T}{HE STATE OF RELIGION} in the West is very uncertain. Opinions abound. Some believe that
Christianity’s influence upon culture is a thing of the past. Many describe the West as
“post-Christian.” Still others seem to think that the Evangelical church has won, in the sense in
that people today would identify Evangelical beliefs as “Christian,” as distinguished from
liberalism. Others, fully informed about declining numbers, sense that North America is poised for
another great work of the Spirit. In terms of the book of Revelation, many in the church are unsure
whether they are the new Israel or the old Babylon. Both directions need to be explored. First,
however, we must attempt to grasp some features of the twentieth century and its story.

As the twentieth century began, conservative Christians struggled with emerging progressive or
liberal Christians who appeared to be capturing the culture. Liberals tapped into a sense of moral
and spiritual progress that awaited flourishing Western civilization. These liberals claimed to
grasp the moral cause and vision of Jesus without being distracted by the distorting doctrine and
liturgy of traditionally minded believers. Other liberals generously recognized the sincere and
genuine religious capacity of every individual person; they reasoned that, while the outward
expression may be different for Buddhists and Baptists, the inner religious experience was the same
at root. These progressive ideas seemed to capture much of the culture of Northern and Western
Europe and North America. Many influential institutions in the West, such as universities, were
characterized by these dispositions or were given to a materialism that rejected spiritual beings.

Conservatives answered the challenge with revival fervor and intellectual muscle. The fundamentalist
movement countered liberal teaching as best they could (though they ironically shared many pivotal
notions of modernity).

The most effective answer seemed to rest in preaching. Some estimate that a million people responded
to the evangelistic call of Billy Sunday. The most effective evangelist, Billy Graham, would share
the gospel internationally. The most effective apologetic voice, C. S. Lewis, came from the
university classroom to articulate a rendering of traditional Christian beliefs. The most
influential theologian was Karl Barth. While not a conservative evangelical, he served their cause
effectively by calling the church and academy back to revelation in light of the inhumanity of
manmade religion witnessed in the two World Wars.

By the time of C. S. Lewis and Billy Graham, the landscape had changed. Two World Wars displayed a
pernicious evil that weighed heavily against the optimism of the progressives. For a time,
Evangelicals seemed to eclipse their liberal challengers with some sense of confidence that the
conservative message had prevailed. After the wars, it became apparent that Evangelicals were facing
other enemies, contributing to a great fading or passing of vibrant Christian faith. Theorists
attribute a variety of contributing factors to this sense of decline. Some say that an endless
preoccupation with self-indulgent individual prerogative and fulfillment has eroded a sense of
community and moral values. Affluence has given an ever-increasing standard to be chased. Many have
noted that “haves” and “have-nots” alike embrace this endless quest to acquire goods, a failure of
virtue called consumerism. With some exceptions, churches bless this economic quest and avoid the
teaching of Jesus concerning money.

Technology plays a part in the daily lives of Westerners that would be hard to imagine just a few
years ago; the twentieth century saw the emergence of television and the amazing network of material
accessible via the computer. The iPhone was introduced in 2007 and has already become standard
equipment for most college students. Smart phones provide individual, private, portable access to
entertainment, information, and instant social feedback or connection with other viewers.

Technology is part of a globalizing tendency that is changing business and creating a common
culture, especially for young people. Where technology will go next is unknown, and its consequences
are also mysterious. Some see it as democratizing knowledge: it makes information available to more
persons, just as the alphabet printing press did. Evidence leads others to conclude that modern
communications technology causes people to suffer from weakened capacity to concentrate and
diminished creativity. With capacity for good or bad, it may just epitomize and aggravate the
self-centered culture we have already noted. Others say the church is to blame. It seems often only
a reflection of its culture and not an agent of change. Secularism is one of the numerous names and
descriptions of this fading of faith.

\section{SECULARIZATION}

The modern understanding of secularization has a complicated history. A central notion for today’s
understanding is that one can clearly distinguish what is religious from what is non-religious. This
may seem odd to readers who are acclimated to this idea, but at one time pursuing the knowledge of
God was thought to be an undertaking that required an entire life venture, drawing upon any field of
knowledge. Knowing God required a distinctive lifestyle, sacred practices, and a new community.
Today the knowledge of God is thought to be an intellectual or academic venture that requires
neither conversion nor even belief in God! Some would even claim that religious commitments
disqualify or hinder the intellectual pursuit of the idea of God. Scholars blame various movements
for this change: Medieval Scholasticism or its collapse (picturing God as merely one influence among
the many), Protestant Scholasticism (relying upon arid intellectual gymnastics), and the emergence
of the nation-state (truncating the idea of God and the arena where he may be served), to name a few
suspects. But perhaps each of these contributes to today’s circumstances where we see religion as a
distinct body of knowledge to go along with biology or physics, and belonging to a distinctly
private or personal realm.

In today’s world, secularism refers to movements that wish to eliminate or restrict the influence of
religion. The secular and religious are inversely proportioned; when religion is gaining ground,
secularization is losing ground, and vice versa. In this climate, secularly minded critics paint a
scenario rooted in the Enlightenment; as more of the world is accounted for by natural explanations,
people will eventually realize that they do not need a category for the supernatural. Hostile voices
against religion see faith as not merely unnecessary but evil; these angry voices claim that
religion is the root cause for most of the problems the world has faced. To be sure, the church has
faced smarter atheists; but the hostility of today’s atheists and secularists calls for the removal
of religion and even the removal of Christianity from the public record. The description of western
culture as “christophobic” is not a serious exaggeration.

\begin{sidebar}{Smarter Atheists}
David Bentley Hart, in Atheist Delusions: The Christian Revolution and its Fashionable Enemies,
observes that recent harsh critics offer an unfitting reading of the past. They blame religion in
general and Christianity in particular for virtually every evil known to man. Nietzsche (1844–1900)
was wiser and better trained. What was mad about Nietzsche’s “madman” was not that he denied that
there was a God. Nietzsche thought that most educated people had ceased to believe in God. Culture
judged the “God-denier” because it was unable to come to grips with the consequence of doing without
God.

Nietzsche knew that Western civilization’s entire moral vision was rooted in a particular notion of
God and that the morality of the West would need to be entirely reworked. Because Nietzsche knew the
history of the West, his conclusion points in the opposite direction from today’s critics. Belief in
God is not simply the problem with the great Western vision of morality. God is the foundation
without which Western values are unsustainable.

Alister McGrath thinks the hostile rhetoric comes from panic: Christianity is growing, and an
unprecedented number of world-class scientists are coming to affirm God. These facts counter the
prevailing secular mythos, that education will lead people away from God.
\end{sidebar}

Perhaps as disturbing as the hostility to faith is an indifference to Christianity. The culture at
large seems not to notice or care; the language of the Bible that once shaped discourse of all
educated citizens seems missing from everyday life. Numerous studies show that believers have less
knowledge and commitment to key Christian teachings than ever before. Those who still defend the
doctrines or teachings of Christianity ironically offer secular explanations, and co-opt
Christianity into a God-sponsored version of self-help therapy. One tangible expression of
secularization is the decline in church attendance and affiliation. Europe typically reports church
attendance under 10 percent, and in some places under 5 percent. This decline is more remarkable
given that the largest churches in many European capitals are led and attended by immigrants from
the Global South.

\section{SHIFT IN WESTERN CHRISTIANS AND CHURCHES}

The consequences of secularism’s rise held dramatic changes for the church in the West. Among the
privileged Protestant mainline churches—the Episcopalians, Methodists, Presbyterians, United Church
of Christ, major Lutheran bodies, and Christian Church (Disciples of Christ)—the impact was largely
negative. They were perceived as symbols of the old established but irrelevant past. As a result,
the mainline bodies saw their membership figures and their finances decline from the mid-1960s
onward, suggesting that these churches were moving from the mainline to the sideline.

In sharp contrast, conservative Protestants were generally flourishing, including Evangelical,
fundamentalist, charismatic, and Pentecostal church groups. To take one striking example during this
age, the Pentecostal Assemblies of God multiplied from a half million to more than four million. By
1990, figures on Roman Catholicism in the United States had also jumped, in their case to fifty-five
million. Large numbers of these new Catholics could be traced to the explosive growth of the U.S.
Latino population. The Roman Catholic Church in America would grow beyond one-third Hispanic.
Conservative churches explained that their stable and growing numbers were due to their fidelity to
faith and Scripture; others included the notion that conservative churches gave solid answers and
were efficiently organized.

\section{EVANGELICAL RESPONSES: RESCUE, RECLAIMING, AND RELEVANCE}

Collectively, conservative churches exhibited three responses to the encroaching secularism. One
response was rooted in a speculative anticipation of the end of time, or at least the beginning of
the end times. In this popular system, the contemporary world was seen as increasingly lost and
having little to do with God’s plan for the church. Because of the increasing evil and moral
failure, they believed that the removal of the church, or Rapture, would happen very soon. Hal
Lindsey published a string of runaway best sellers, including The Late, Great Planet Earth and Satan
is Alive and Well On Planet Earth. These books created a devout expectation that Jesus must be
coming anytime; this heightened awareness showed that redeeming the fallen world was futile.

Ministry and mission alike must be given to preparing souls before the end comes. Even some
Christian critics thought this evangelical fervor took on Gnostic shades, in that it depreciated
God’s creation and his intention to redeem it; the body as well as the soul was to be redeemed.
After several decades of missed deadlines for Jesus’ return, the movement lost steam, though any
number of headlines involving the Middle East or the European Union could awaken this mindset and
give opportunity for another book to update the situation. Evangelicals did not abandon this
approach to the future, but they generally came to sense that they had at least applied it in a
dubious fashion. Perhaps inconsistently, some of the same personalities participated in the next
approach, seeking to contest Christian claims upon the world and nation. The approaches blend
together in practice.

The next approach was to battle secular forces on political and cultural turfs. Jerry Falwell was
pastor of the Thomas Road Baptist Church in Lynchburg, Virginia, but after creating a political
action group called the Moral Majority in 1979, he became the first prominent spokesman for the
so-called Religious Right. The label stood for a coalition of fundamentalist, Pentecostal,
Evangelical, and Catholic Christians, driven by concern for the decline in American morality, who
had become extremely active in the political arena. The core of the movement was a loose alliance of
groups led by the Moral Majority. Clustered around an agenda defending traditional moral values and
conservative political goals were The Christian Voice, led by Robert Grant; Concerned Women for
America, under the leadership of Beverly LaHaye; and the Freedom Council, formed by Pat Robertson, a
televangelist who became an unsuccessful candidate for the Republican presidential nomination in
1988.

The passion of the Religious Right lay in their perception that the United States was falling under
the influence of secular humanism and that traditional family values were under attack in the media
and the public schools. Several explosive national issues seemed to ignite this new conservative
reaction. First, in 1973 the U.S. Supreme Court’s decision in the Roe v. Wade case agreed with Jane
Roe, a young single Texas woman, that her right to privacy included her right to terminate her
pregnancy by abortion.

Many Catholics and Protestants, who held that human life begins at conception, were shocked and
dismayed by the decision. Second, in 1978 a ballot proposition in California tried to expand legal
protection of homosexuals. When a group of conservative pastors organized to defeat the measure, the
Internal Revenue Service warned them that the tax-exempt status of their churches was endangered by
their political activity. The pastors thought this was reason enough to form The Christian Voice.
They were unwilling to give up their fight. Finally came the battle over the Equal Rights Amendment,
the movement to grant women the legal protection of an amendment to the Constitution. Like many
conservative Christians, Jerry Falwell was opposed to the wording of the amendment adopted by
Congress. He felt that it would allow homosexual marriages and adoptions and would allow women to be
drafted into the armed services. So in Virginia he fought successfully against the ratification of
the amendment.

These three issues—abortion, homosexual rights, and feminism—served to mobilize conservative
Christians for battles in the political arena, soon to expand to include other moral and political
issues. To achieve their goals, the Religious Right relied heavily on the ministries of
televangelists. Falwell’s Old Time Gospel Hour and Robertson’s Christian Broadcasting Network (CBN)
and 700 Club led a host of Christian-sponsored radio and television outlets in promoting the
conservative moral and political agenda. On radio, a child psychologist-turned-broadcaster, Dr.
James Dobson, enlarged his radio ministry called Focus on the Family into a powerful voice for
traditional family values. Perhaps most significant politically, these groups succeeded in educating
and mobilizing fundamentalists and Pentecostals, a segment of the American population that had once
been politically inactive.

A third approach involves another change of course or strategy. Numerous Evangelicals sought to be
less confrontational and would attempt to engage the new changing culture on its own terms. These
churchmen observed that they ministered within the age of individual self-expression. They created
user-friendly churches that made religion almost totally a matter of personal choice. Desiring to
evangelize, practitioners performed market analysis in an effort to offer an experience that would
meet with the approval of the target audience. Some fashioned an experience that would be
comfortable to a religious seeker. A great variety of churches adopted a practical self-help
approach; a special version of allegory was embraced to extract techniques from the Bible, which
suddenly became a manual for success in life.

The megachurch phenomenon illustrates the third approach of seeking relevance. With the decline of
denominations in American public life and the increasing privatization of religion, large churches
gained an increasing share of churchgoing America, or, as some crassly called it, “the religious
market.” These large churches grew at least in part because they shed the negative image of
denominational Christianity and appealed to popular religious tastes. Like the seventy-six million
members of the post-World War II generation, megachurches liked to think of themselves as
independent and highly individualized.

With attenders in the thousands, Sunday morning services in these churches were usually well
attended. But their buildings were filled the rest of the week, too, with Bible classes, support
groups, field trips for seniors, weight loss classes, and children’s activities. The appeal to
popular taste was revealed in several common characteristics: First, these congregations seldom
carried a denominational label. They much preferred “chapel,” “center,” or “community” on the sign
out front. The name was a symbol of their openness to people with diverse backgrounds and problems:
divorce, addictions, and depression. Second, the worship in these large congregations was marked by
fast paced and enthusiastic popular religious music. The choice of music has been increasingly
influenced by a Christian commercial music industry. Third, they were built around the attractive
ministry of a magnetic preacher who possessed a winsome personality. The sermons stressed the
Bible’s application to day-to-day life. Fourth, these large churches give an image of efficiency and
affluence. Buildings were often new; the staff of ministers was well trained and effective; services
were available for every imaginable need.

Historians tend to view this development as another example of Americans’ privatization of faith.
Even in the excitement of a large crowd, attenders of a megachurch were looking for faith that
served the private life: help on child rearing, family unity, and personal emotions. Experienced
churchgoers found these churches uncomplicated, by comparison with the small neighborhood churches
where they had formerly attended. Americans sought a private house, a private means of
transportation, a private garden, a private laundry, and self-service stores. Even within families,
Americans had come to expect that each member of the family should have a separate room and even a
separate television, phone, and car, at least when economically possible.

Many megachurches continue to flourish, but many others have plateaued. A new generation may well
seek community to compensate for individualism and isolation; some seek a new approach that will
appeal to the emerging generation. Hundreds of new churches also appeared throughout the United
Kingdom and the United States, all attempting to reach out to the postmodern generation. “Those who
fail to change their ministry ways,” said the new voices, “risk hiding the gospel behind forms of
thought and modes of expression that no longer communicate with the new emerging generation.”

Emerging churches were intent upon emphasizing feelings and affections rather than rationality and
linear thought, on personal experience over propositional truths, on inclusion rather than
exclusion, and on participation in corporate worship in contrast to lost-in-the-crowd, megachurch
individualism.

\section{TROUBLING SIGNS}

Sobering signs show that conservative evangelical efforts (to prepare for spiritual rescue from a
dying world or to reclaim a Christian character for the nation by political victory or to connect
with the lost culture) have not stemmed the tide of secularization. Taken together, numerous signs
indicate that “old-time religion” may be moving swiftly toward a minority status within the culture.
Conservatives could no longer look at church decline as a mainline or liberal church problem. When
many conservative denominations began to decline, haunting questions emerged: Will conservative
denominations follow liberal church declines? Will the North American Christian landscape soon
resemble that of Europe? Conservatives could no longer take solace in slowing declines in the early
1990s. Recent Pew Research Center numbers measure religious affiliation at its lowest since their
earliest records in 1930. Now one in five Americans claims to have no religious affiliation; for
adults under age thirty, the number jumps to one in three. Also conservative Christians have lost
ground on the three initiatives targeted by the Religious Right: abortion, homosexual rights, and
feminism. Legal limitations upon same-sex marriage received serious damage in June, 2013.
California’s Proposition Eight and the Defense of Marriage Act were both set aside in Supreme Court
action. Additionally, conservative believers in the West feel threatened by several ongoing issues.

\section{ONGOING STRUGGLES}

The institutions of the family and marriage have suffered. Traditional notions of couples being
married before they live together and the enduring nature of marriage are not recognized by the
larger culture or Christian church members in a consistent fashion. The general culture abandoned
the traditional idea of marriage long before it came under siege by same-sex interests.
Single-parent families (which have reached one in five) and working women (at least forty percent of
the work force) add to the destabilized nature of the home.

On the morning of September 11, 2001, Americans were dramatically confronted with radical terrorism.
They discovered that religion, even religious zealots, would shape the political landscape for the
foreseeable future. Two airliners filled with jet fuel crashed into the two towers of New York City
World Trade Center. Having failed to bring down the buildings in 1993, this time they succeeded. A
third plane flew into the Pentagon, and passengers struggled with hijackers before a fourth plane,
intended for the White House, crashed in southwest Pennsylvania. Almost every American can see the
Trade Center buildings collapse in their mind’s eye.

Within days of 9/11, President George W. Bush declared war on the terrorists and, after rallying a
coalition of largely Western nations, the president—less than a year in office—sent troops to
Afghanistan, where the terrorists called the Taliban had found a haven and a training ground. With
the cooperation of key Afghan warlords, the American-led coalition was able to topple the Taliban
and drive it from power within a matter of months. Osama bin Laden, leader of the most feared
jihadists, was later found and shot, but only after he had inspired and directed new attacks upon
the “infidels.”

Within months, a new government was in place in Kabul, the Afghan capital. The invasion of Iraq, the
reported haven of other jihadists, and the toppling of Saddam Hussein from power soon followed in
the American-led coalition’s war against terror. The war in Iraq, however, a much larger and
stronger country, proved to be far more complex and difficult than the Afghan conflict.

\begin{sidebar}{Modernity and Islam}
One of the common explanations for the term modernity is religious and political in nature. In this
approach, a nation-state is modern if it tolerates differing religious expressions under its larger
national umbrella. The state is to some degree neutral or generic in a religious sense and allows
differing voices, such as Islam and Christianity to live in the same state. This concept is helpful
in understanding Islam. After World War I, almost the entire Arab world was placed under a state
(some very contrived, placing antagonistic groups together under one flag) that was overseen by a
European nation (and in the eyes of most Arabs, these are also Christian States). This arrangement
is part of a larger picture called imperialism.

While Christians had altered their practice to agree to a modern state not long after the
Reformation (the Thirty Years’ War), Islam is by nature almost necessarily social and political;
Islam’s voices typically think that Islam must be translated into government-sanctioned laws to be
faithful. For a time, the Muslim world lived in modern states sponsored by European countries at
gunpoint. Rulers such as Sadat in Egypt and the Shah of Iran brutally suppressed radical voices to
make these modern states possible. Egypt and Iran were thought to be westernized. Once the Arab
peoples began to exercise political independence, radical voices within Islam proved effective in
seizing control. Iran’s revolution brought a sudden change from a westernized nation to one governed
by strict Islamists. Time will tell which countries among the Arab people will select strict Islamic
states or embrace some sort of modern solution.

At present, the voices of radical Shiite Muslims seem to shape much of the political landscape and
the culture. Modern Muslims, which most students view as the majority, seem quite stifled. The rise
of radical Islam has changed the religious landscape across the globe, and many Arab countries are
now radicalized. It should be observed that Islamic rulers have historically tolerated Christian
citizens as long as they submitted to being second class citizens within Muslim nations.
\end{sidebar}

Today, Christians in the Arab world are suffering. In 1900, the Middle East and North Africa were 15
percent Christian; one hundred years later, estimates were 5 percent. The once-significant numbers
of Arab Christians have declined as they have been persecuted into submission, exiled, or murdered.
Perhaps as significant is the presence of Muslim peoples living in the West with missionary and even
militant purpose. Christians are wise to be alert to the risk and danger that radical Islam holds
for them in the Arab or Western worlds. Christians will face a still greater danger of infidelity to
Christ if Muslims do not encounter the love of Christ in tangible witness and care.

Persistent testimony of Muslim conversion and openness to the gospel is difficult to measure;
reports of Muslims encountering Jesus in their dreams, the curiously positive reception some that
evangelical Christians receive in the Arab world, the persistence of underground churches, and
Christians willing to serve and live among Muslims give some encouragement for a continued Christian
presence in the region. These Christians who are called to live in the Arab world are some of the
champions of the missionary venture.

Added to the awareness and danger of Muslim radicals, the church was also called to face a new wave
of immigration to the country. The Hart-Cellar Act was a part of the Civil Rights legislation. The
Act stimulated immigration from Asia, the Middle East, and Africa. Illegal immigration from South
and Central America also accelerated racial diversity. Though voices from the church led the way to
progress in the matter of race, the evangelical church has largely maintained racial segregation as
well as segregation by economic standing. Race and economic class go hand in hand; for example, a
black family will be more readily accepted in a largely Anglo church if they share a similar
lifestyle.

Also as missions were exploding around the globe, Western Christians were beginning to struggle with
the practice. Evangelical denominations and mission-sending organizations were facing budget cuts.
Evangelicals wondered if missionaries should address more humanitarian concerns over concerns about
salvation; they also struggled as to what degree missionaries imposed their cultural bias upon the
nationals; they also were humbled by more extensive encounters with world religions. Racial and
geographical diversity, the awareness of Islam, and globalizing technology all bring these issues
home. Missionary zeal seemed easier overseas; it seemed more complicated when the world had come
near geographically and grown smaller technologically.

\section{DEFEATED OR RECONSTITUTING?}

News of Christianity’s demise in the West may be premature. Minimally, it will be influential
because of its money and educational institutions. Western degrees are often still privileged in the
Global South. More complete assessment is difficult, though several observations may help. The
cultural environment is increasingly secular and hostile to faith. Christians looking for
Christendom (a Christian domain marked by an organic and cooperative partnership between the state
and church) will be disappointed. Christians and their convictions can expect to be more
marginalized and persecuted. Pastors being prosecuted for hate speech and Christian schools losing
federal dollars seem almost unavoidable. The mainline denominations will experience general decline,
although they will vary regarding their acclimation to cultural norms; some may return instinctively
to the great tradition of the church and rediscover the gospel (this approach is called
post-liberalism).

Certain churches in the Pentecostal tradition, nondenominational circles, and conservative
denominations will continue to grow. The real progress will rest in new church starts, which are
grossly under-reported. These churches are pliable and portable; they can meet in store fronts and
in cell groups. They are almost impossible to track or count.

A personal parable may help illustrate. I (R. L. Hatchett) was raised in a neighborhood Baptist
church. Our church, the Methodists, and Church of Christ churches were the norm in our humble
circumstance. Our church buildings looked like churches and were nicer than our own homes. When a
charismatic nondenominational church began to meet in an old grocery store space, the contrast was
unmistakable. They had a long sign with too many words. They were even poorer than we were. They
were unconnected to larger bodies and institutions (like mission programs and seminaries). Most of
all, they were the periphery or exception, whereas we were the norm. Now almost fifty years later,
these upstart churches are not the periphery but the numerical norm. The immediate future belongs to
groups like this that are independent and entrepreneurial.

The new norm is composed of the Pentecostals and non-denominationals; their emergence signals a
change in the understanding of Christianity for the future in the West. We are engaged in a
reconstruction period that follows this sociological revolution; the evangelical population will now
be found predominantly in this mobile, minority-rich matrix. They display a certain Anabaptist-like
character. They are not all pacifists (though some forget that many Pentecostals were conscientious
objectors in World War I) and some may nostalgically look for the good days of a Christian America,
but they see themselves as countercultural; they focus on discipleship and community; they will
partner to meet needs but do not have the political disposition and orientation of older
Evangelicals.

Other signs of life are found in a variety of places. The rigorously academic and spiritually devout
young Catholics in America seem strong in mind and heart and number. One is encouraged as well by
youth gatherings such as Passion. In 2013, sixty-four thousand college age students met in Atlanta
during their Christmas break. Each Passion conference, whether in the states or overseas, fills to
capacity. Numerous initiatives call for renewal and revival. An ecumenical group, Renovaré, crosses
denominational lines to learn about spiritual formation. We pray for the Spirit to move.

\readinghead{Suggestions for Further Reading}
\begin{readinglist}
\entry Bellah, Robert N. et al. Habits of the Heart: Individualism and Community in American Life.
  New York: Harper \& Row, 1985.
\entry Bennett, William J. The Index of Leading Cultural Indicators. New York: Simon and Schuster,
  1994.
\entry Collins, Kenneth J. The Evangelical Moment: A Promise of an American Religion. Grand Rapids:
  Baker Academic, 2005.
\entry George, Timothy, ed. Pilgrims on the Sawdust Trail: Evangelical Ecumenism and the Quest for
  Christian Identity. Grand Rapids: Baker, 2004.
\entry *Noll, Mark. The Old Religion in a New World: the History of North American Christianity.
  Grand Rapids: Eerdmans, 2002.
\entry Shelley, Bruce, and Marshall Shelley. The Consumer Church. Downers Grove, IL: InterVarsity
  Press, 1992.
\entry *Weigel, Georg. The Cube and the Cathedral: Europe, America, and Politics Without God. New
  York: Basic Books, 2005.
\entry Wuthnow, Robert. The Struggle for America’s Soul, Evangelicals, Liberals, and Secularism.
  Grand Rapids: Eerdmans, 1989.
\end{readinglist}
